% ===========================================================================
% Shared CONTENT preamble.
%
% This holds only the packages and macros that the *content* depends on --
% math, pseudocode styling, the syntax-highlight colour palette, and the
% semantic macros (\code, \bl, \LineComment). It carries NO document-level
% typesetting (geometry, hyperref, biblatex, title, doclicense, fancyhdr).
%
% It is \input by two places, which is the whole point of factoring it out:
%   1. cpsat_search_core.tex        -> the print (PDF) build
%   2. md_build/ standalone builder -> renders each pseudocode/figure float to
%                                       an image, guaranteed identical to print.
%
% Keep this file in sync with the source: the macro definitions below are
% copied verbatim from the original cpsat_search_core.tex preamble.
% ===========================================================================

% --- Fonts (fontspec defaults = Latin Modern, matching the article) --------
\usepackage{fontspec}
\usepackage{microtype}

% --- Math ------------------------------------------------------------------
\usepackage{amsmath}
\usepackage{amssymb}

% --- Pseudocode ------------------------------------------------------------
\usepackage{algorithm}
\usepackage{algpseudocode}

% pseudocode styling: tighter indent, comment style
\algrenewcommand\algorithmicindent{1.1em}
\algrenewcommand\algorithmiccomment[1]{\hfill\textit{\footnotesize $\triangleright$\,#1}}
% Phase comment: a normal numbered, indented code line (like in an editor), set apart
% only by style --- a grey italic /* ... */ block reusing the listings' comment colour.
\algnewcommand{\LineComment}[1]{\State\textcolor{codecom}{\textit{\footnotesize /*~#1~*/}}}

% --- Colour palette (shared with the code/listings styling) ----------------
\usepackage{xcolor}
\definecolor{codekw}{HTML}{0033B3}      % keywords (blue)
\definecolor{codebuiltin}{HTML}{6A2570} % builtins (purple)
\definecolor{codestr}{HTML}{B35C00}     % strings (orange)
\definecolor{codecom}{HTML}{6E737A}     % comments (grey)
\definecolor{codenum}{HTML}{1E7C50}     % numbers (green)
\definecolor{codedecorator}{HTML}{735A12}% decorators

% --- Semantic inline macros ------------------------------------------------
\usepackage{seqsplit}
% inline monospace helper; long identifiers made breakable via \seqsplit.
\DeclareRobustCommand{\code}[1]{\texttt{\seqsplit{#1}}}

% bound-literal brackets: \bl{x \ge 5} -> [[x >= 5]]
\newcommand{\bl}[1]{[\![#1]\!]}
